\begin{proposicion}[Ratio of normalizers]\label{coste-normalizadores}
Under the same hypotheses,
\[
 \log(\Gamma_4/\Gamma_2)=O_{C,\eta,\gamma}(\varepsilon_{\rm out}+k^{-1}).
\]
\end{proposicion}
\begin{proof}
Write $a_2=\sum_e p_e^2$. We first compare the row and column blocks. On the orthogonal complement
of their common kernel, the derivative of $\log\det$ is the trace of the inverse times the
derivative of the matrix. Lemma~\ref{coste-inversa} and $\tau_4(p)-\tau_2(p)=2p^2+O(p^3)$ give
\[
 \begin{split}
 \log\frac{\pdet S_{R,4}}{\pdet S_{R,2}}
 &=\int_0^1\sum_e(\tau_4(p_e)-\tau_2(p_e))\,
                   \iota_e^TS(s)^+\iota_e\,ds\\
 &=4a_2+O(k^{-1}).
 \end{split}
\]
Indeed, $\iota_e^TS(s)^+\iota_e=2+O(k^{-1})$, $a_2=O(1)$, and
$\sum_e p_e^3\le p_{\max}a_2=O(k^{-1})$.

Write $\tau_e=\tau_4(p_e)$, $G=S_{R,4}^+$, and
\[
 V=A_R^T\operatorname{diag}(\tau)D,\qquad
 S_D=\operatorname{diag}(h)+D^T\operatorname{diag}(\tau)D,\qquad
 U=S_D-V^TGV .
\]
Each $\iota_e$ is orthogonal to the zero mode, and the columns of $V$ are linear combinations of
these vectors. In an orthonormal basis for the orthogonal complement of the zero mode, block
elimination gives $\pdet S_4=(\pdet S_{R,4})\det U$. For any vector $z$ indexed by diagonals,
the form $z^TU z$ is the minimum of the quadratic form of $S_4$ at $(y,z)$ over $y$ orthogonal
to the zero mode. The minimizer is $-GVz$. The lower bound of Lemma~\ref{coste-estatico}
therefore gives $U\ge cI$ with the same $c>0$.

Put $A=S_D-I$ and $T=V^TGV$, so that $U=I+A-T$. On the diagonal,
$A_{ff}=\sum_{e\in f}(\tau_e-p_e)$. Off the diagonal, only the entries between the two diagonal
families can be nonzero. Such an entry equals $\tau_e$ when the lines meet in an available square
$e$ and is zero otherwise. It follows that
\[
 \|A\|_{\rm op}=O(\varepsilon_{\rm out}),\qquad
 \operatorname{tr}A=6a_2+O(k^{-1}),\qquad
 \operatorname{tr}A^2=2a_2+O(k^{-1}).
\]
Indeed, the sum of absolute values in each row is $O(\ell_f)$, and the trace counts
$\tau_e-p_e=3p_e^2+O(p_e^3)$ twice per square. For the trace of the square, the off-diagonal
entries contribute $2\sum_e\tau_e^2=2a_2+O(k^{-1})$. The diagonal contribution is $O(k^{-1})$
because $\max_f A_{ff}=O(k^{-1})$ and $\sum_f A_{ff}=O(1)$.

Each nonzero entry of $V$ equals $\tau_e$, since a row or column and a diagonal determine at most
one square. The row sums of $V$ are at most $2\max_u\sum_{e\in u}\tau_e=O(1)$, where $u$ ranges
over rows and columns. The column sums are at most $2\max_f\sum_{e\in f}\tau_e=O(\varepsilon_{\rm out})$.
The bound on the spectral norm in terms of the row and column sum norms gives
$\|V\|_{\rm op}^2=O(\varepsilon_{\rm out})$. Moreover,
$\|V\|_{\rm HS}^2=4\sum_e\tau_e^2=O(1)$. Since $G$ is positive semidefinite with norm at most
$\gamma^{-1}$,
\[
 T\ge0,\qquad \|T\|_{\rm op}=O(\varepsilon_{\rm out}),\qquad
 \operatorname{tr}T=O(1),\qquad
 \operatorname{tr}T^2\le\|T\|_{\rm op}\operatorname{tr}T
             =O(\varepsilon_{\rm out}) .
\]

To find the leading term in the trace of $T$, we use the incidence structure.
By Lemma~\ref{coste-inversa},
\[
 \operatorname{tr}T
 =\sum_f\sum_{e,e'\in f}\tau_e\tau_{e'}\iota_e^TG\iota_{e'}
 =4a_2+O(\varepsilon_{\rm out}+k^{-1}) .
\]
The pairs $e=e'$ contribute $2\sum_e\tau_e^2(2+O(k^{-1}))$ because each square lies on two
diagonals. Distinct squares on the same diagonal have distinct rows and columns, so their mixed
term is $O(k^{-1})$. The total absolute error from these pairs is bounded by
\[
 \frac{C'}k\sum_f\Bigl(\sum_{e\in f}\tau_e\Bigr)^2
 \le\frac{C''}k\sum_f\ell_f^2
 \le\frac{C''}k\,\varepsilon_{\rm out}\sum_f\ell_f
 =O(\varepsilon_{\rm out}),
\]
since $\sum_f\ell_f=2k$.

For $E=A-T$ we have $\|E\|_{\rm op}=O(\varepsilon_{\rm out})$ and $\operatorname{tr}E^2=O(1)$.
Also $|\operatorname{tr}(AT)|\le\|A\|_{\rm op}\operatorname{tr}T=O(\varepsilon_{\rm out})$
by positivity of $T$. Consequently,
\[
 \operatorname{tr}E=2a_2+O(\varepsilon_{\rm out}+k^{-1}),\qquad
 \operatorname{tr}E^2=2a_2+O(\varepsilon_{\rm out}+k^{-1}).
\]
If $\xi$ is an eigenvalue of $E$, then $1+\xi\ge c$. Taylor's formula gives
$|\log(1+\xi)-\xi+\xi^2/2|\le|\xi|^3/(3\min\{c,1\}^3)$.
Upon summing, the error is $O(\|E\|_{\rm op}\operatorname{tr}E^2)=O(\varepsilon_{\rm out})$.
In particular,
\[
 \log\det U=a_2+O(\varepsilon_{\rm out}+k^{-1}).
\]
Finally, the quadratic coefficient of $\log(1-dp)-d\log(1-p)$ is $-d(d-1)/2$, with an $O(p^3)$
remainder for $d=2,4$. Thus $C_4-C_2=-5a_2+O(k^{-1})$. The factor $\sqrt{4\pi k}$ is common
to $\Gamma_4$ and $\Gamma_2$. Combining the three terms gives
\[
 \log(\Gamma_4/\Gamma_2)
 =-\tfrac12\bigl((-5+4+1)a_2\bigr)
      +O(\varepsilon_{\rm out}+k^{-1}),
\]
as claimed.
\end{proof}
